\chapter{Mínimos Quadrados Ordinários}
\label{cap:ols}

\section{Sintaxe}

\begin{lstlisting}
ols(fórmula, df)
ols(fórmula, df, cov=tipo)
ols(fórmula, df, cov=cluster, cluster=coluna)
ols(fórmula, df, if = condição)
\end{lstlisting}

Alias: \lstinline{reg()} é equivalente a \lstinline{ols()}.

\section{Fórmula}

A fórmula segue a notação padrão \texttt{y \textasciitilde{} x1 + x2 + ...}:

\begin{lstlisting}[caption={OLS básico}]
input df
Y X
10 2
12 3
8  1
15 5
11 2
14 4
end

ols(Y ~ X, df)
\end{lstlisting}

A constante é incluída automaticamente. Para múltiplas variáveis:

\begin{lstlisting}
ols(Y ~ X1 + X2 + X3, df)
\end{lstlisting}

A fórmula pode ser passada como string:

\begin{lstlisting}
let formula = "Y ~ X1 + X2"
ols(formula, df)
\end{lstlisting}

\section{Erros padrão}

O tipo de erro padrão é controlado pela opção \lstinline{cov=}:

\begin{center}
\begin{tabular}{ll}
\toprule
\textbf{Opção} & \textbf{Erro padrão} \\
\midrule
(padrão)              & OLS clássico (homocedasticidade) \\
\texttt{cov=robust}   & Robusto à heterocedasticidade (HC1) \\
\texttt{cov=HC1}      & Idem — alias explícito \\
\texttt{cov=HC3}      & HC3 (mais conservador em amostras pequenas) \\
\texttt{cov=cluster, cluster=col}  & Clusterizado por \texttt{col} \\
\bottomrule
\end{tabular}
\end{center}

\begin{lstlisting}[caption={Variantes de erro padrão}]
ols(Y ~ X, df)                           // clássico
ols(Y ~ X, df, cov=robust)              // HC1
ols(Y ~ X, df, cov=HC3)                 // HC3
ols(Y ~ X, df, cov=cluster, cluster=id) // clusterizado
\end{lstlisting}

\section{Amostra condicional}

A opção \lstinline{if =} restringe a estimação a um subconjunto sem modificar o DataFrame:

\begin{lstlisting}
ols(Y ~ X, df, if = grupo == 1)
\end{lstlisting}

\section{Capturando o resultado}

O resultado de \lstinline{ols()} é um objeto que pode ser atribuído a uma variável
e usado em pós-estimação:

\begin{lstlisting}
let m = ols(Y ~ X, df)
\end{lstlisting}

Quando atribuído, \lstinline{ols()} não imprime a tabela — apenas armazena o modelo.
Para exibir o resultado explicitamente:

\begin{lstlisting}
let m = ols(Y ~ X, df)
display m
\end{lstlisting}

\section{Pós-estimação}

\subsection{\texttt{predict}}

Extrai valores ajustados ou resíduos do modelo:

\begin{lstlisting}
let m = ols(Y ~ X, df)

let yhat = predict(m, "xb")        // valores ajustados
let ehat = predict(m, "residuals") // resíduos
\end{lstlisting}

Aliases aceitos: \texttt{"fitted"} para valores ajustados; \texttt{"resid"} ou
\texttt{"e"} para resíduos.

\subsection{\texttt{vif}}

Fator de inflação da variância — diagnóstico de multicolinearidade:

\begin{lstlisting}
let m = ols(Y ~ X1 + X2 + X3, df)
vif(m)
\end{lstlisting}

\subsection{\texttt{test}}

Teste de hipótese linear sobre os coeficientes (teste F):

\begin{lstlisting}
let m = ols(Y ~ X1 + X2 + X3, df)
test(m, X2 = 0, X3 = 0)    // H0: X2 = X3 = 0
\end{lstlisting}

\subsection{\texttt{lincom}}

Combinação linear de coeficientes com erro padrão:

\begin{lstlisting}
let m = ols(Y ~ X1 + X2, df)
lincom(m, X1 + X2)
\end{lstlisting}

\subsection{\texttt{margins}}

Efeitos marginais médios (para modelos não-lineares; no OLS equivale aos coeficientes):

\begin{lstlisting}
let m = ols(Y ~ X, df)
margins(m)
\end{lstlisting}

\section{Tabela de resultados com \texttt{esttab}}

\lstinline{esttab} exibe múltiplos modelos em uma tabela comparativa:

\begin{lstlisting}[caption={Comparação de especificações}]
let m1 = ols(Y ~ X, df)
let m2 = ols(Y ~ X, df, cov=robust)
let m3 = ols(Y ~ X1 + X2, df)

esttab(m1, m2, m3)
\end{lstlisting}

Alternativamente, use \lstinline{eststo} para acumular modelos e depois exibir:

\begin{lstlisting}
eststo(ols(Y ~ X, df))
eststo(ols(Y ~ X1 + X2, df))
esttab()
\end{lstlisting}

\section{Bootstrap}

Erros padrão por bootstrap:

\begin{lstlisting}
bootstrap(ols, Y ~ X, df, n=1000)
bootse(m, n=500)
\end{lstlisting}
